%%--------------------------------------------------------------------
%1--------------------------------------------------------------------
\vspace*{\fill}
\begin{glyph}{Copy (写)}{copy}\label{glyph:copy}
Copy, photograph.
\end{glyph}

\begin{comps}
\comprtwocone & 冖 + 与 \\\hline
\comprthreecone & Cover/Crown + Confer (\ref{glyph:confer}) \\\hline
\end{comps}

\begin{remglyph}
\hooglig{A copy} was \remcom{conferred} to me, and it was \remcom{covered} in an envelope.\\\hline
\end{remglyph}

\begin{prons}
\pronsrtwocone & \textbf{シャ (SHA)} & \textbf{うつ (utsu)}\\\hline
\pronsrthreecone & --- & --- \\\hline
\end{prons}

\begin{remprons}
\hooglig{A photograph} of where the \comkun{roots} were \comon{shut} in.\\\hline
\end{remprons}

%{\middel}
% &  \mbox{()}\\\hline
\begin{somme}
写す & \textit{copy; take a photograph} & うつ{\middel}す \mbox{(utsu{\middel}su)}\\\hline
写真 & copy + true = \textit{photograph} & シャ{\middel}シン \mbox{(SHA{\middel}SHIN)}\\\hline
写生 & copy + life = \textit{sketch; sketching} & シャ{\middel}セイ \mbox{(SHA{\middel}SEI)}\\\hline
\end{somme}

%\begin{decomp}{}
%\\\hline
%\\\hline
%\\\hline
%\multicolumn{}{|r|}{\textit{}}\\\hline
%\end{decomp}

\end{multicols}
\vhrulefill{2pt}
%%--------------------------------------------------------------------
%2--------------------------------------------------------------------
\begin{multicols}{2}
\begin{glyph}{Court (廷)}{court}\label{glyph:court}
Court.\\\hline
This 漢字 can refer to courts of law (as in the initial two compounds below), and imperial courts (as in the last).
\end{glyph}

\begin{comps}
\comprtwocone & 廴 + 士 + 丿 \\\hline
\comprthreecone & Long stride + Scholar + Slash \\\hline
\end{comps}

\begin{remglyph}
The court proceeding travesed a long distance according to the scholars slashing each others' throats.\\\hline
\end{remglyph}

\begin{prons}
\pronsrtwocone & \textbf{テイ (TEI)} & --- \\\hline
\pronsrthreecone & --- & --- \\\hline
\end{prons}

\begin{remprons}
\hooglig{The court} proceedings were \comon{taped}.\\\hline
\end{remprons}

%{\middel}
% &  \mbox{()}\\\hline
\begin{somme}
法廷 & law + court = \textit{law court} & ホウ{\middel}テイ \mbox{(H\={O}{\middel}TEI)}\\\hline
出廷 & exit + court = \textit{appear in court} & シュッ{\middel}テイ \mbox{(SHUT{\middel}TEI)}\\\hline
朝廷 & morning (dynasty) + court = \textit{the imperial court} & チョウ{\middel}テイ \mbox{(CH\={O}{\middel}TEI)}\\\hline
\end{somme}

%\begin{decomp}{}
%\\\hline
%\\\hline
%\\\hline
%\multicolumn{}{|r|}{\textit{}}\\\hline
%\end{decomp}

\end{multicols}
\vspace*{\fill}
\pagebreak
%%--------------------------------------------------------------------
%3--------------------------------------------------------------------
\vspace*{\fill}
\begin{multicols}{2}
\begin{glyph}{Sideways (横)}{sideways}\label{glyph:sideways}
Sideways, horizontal.
\end{glyph}

\begin{comps}
\comprtwocone & 木 + 黄 \\\hline
\comprthreecone & Tree + Yellow \\\hline
\end{comps}

\begin{remglyph}
\hooglig{Horizontal} \remcom{yellow} \remcom{tree}.\\\hline
\end{remglyph}

\begin{prons}
\pronsrtwocone & \textbf{オウ (\={O})} & \textbf{よこ (yoko)}\\\hline
\rye{3}{The 訓読み is the more common reading.}
\pronsrthreecone & --- & --- \\\hline
\end{prons}

\begin{remprons}
\hooglig{Horizontal} beaches in \comkun{Yokohama} are \comon{awesome}.\\\hline
\end{remprons}

%{\middel}
% &  \mbox{()}\\\hline
\begin{somme}
横 & \textit{sideways; horizontal} & よこ \mbox{(yoko)}\\\hline
横切る & sideways + cut = \textit{go across; traverse} & よこ　ぎ{\middel}る \mbox{(yoko gi{\middel}ru)}\\\hline
横転 & sideways + roll = \textit{turn sideways; barrel roll} & オウ{\middel}テン \mbox{(\={O}{\middel}TEN)}\\\hline
\end{somme}

%\begin{decomp}{}
%\\\hline
%\\\hline
%\\\hline
%\multicolumn{}{|r|}{\textit{}}\\\hline
%\end{decomp}

\end{multicols}
\vhrulefill{2pt}
%%--------------------------------------------------------------------
%4--------------------------------------------------------------------
\begin{multicols}{2}
\begin{glyph}{Certain (必)}{certain}\label{glyph:certain}
Certain, sure.  Despite its simplicity, writing it is difficult due to 心 being written in an entirely different stroke order.
\end{glyph}

\begin{comps}
\comprtwocone & 心 + 丿 \\\hline
\comprthreecone & Heart + Slash \\\hline
\end{comps}

\begin{remglyph}
\hooglig{Certainly} \remcom{heart} and \remcom{stroke} attacks will do the trick.\\\hline
\end{remglyph}

\begin{prons}
\pronsrtwocone & \textbf{ヒツ (HITSU)} & \textbf{かなら (kanara)}\\\hline
\pronsrthreecone & --- & --- \\\hline
\end{prons}

\begin{remprons}
\hooglig{Certain} \comkun{canna land} \comon{hits}.\\\hline
\end{remprons}

%{\middel}
% &  \mbox{()}\\\hline
\begin{somme}
必ず & \textit{certain; for sure} & かなら{\middel}ず \mbox{(kanara{\middel}zu)}\\\hline
必要 & certain + essential = \textit{essential; necessary} & ヒツ{\middel}ヨウ \mbox{(HITSU{\middel}Y\={O})}\\\hline
必死 & certain + death = \textit{certain death; desperate} & ヒッ{\middel}シ \mbox{(HIS{\middel}SHI)}\\\hline
\end{somme}

%\begin{decomp}{}
%\\\hline
%\\\hline
%\\\hline
%\multicolumn{}{|r|}{\textit{}}\\\hline
%\end{decomp}

\end{multicols}
\vspace*{\fill}
\pagebreak
%%--------------------------------------------------------------------
%5--------------------------------------------------------------------
\vspace*{\fill}
\begin{multicols}{2}
\begin{glyph}{With (以)}{with}\label{glyph:with}
This 漢字 works together with those that follow it, and with its companion sets some type of condition on the 漢字 preceding them.
\end{glyph}

\begin{comps}
\comprtwocone & 丨 + 丶 + 人 \\\hline
\comprthreecone & Pole + Dot + Person/Man \\\hline
\end{comps}

\begin{remglyph}
\hooglig{With} \remcom{jelly-beans} skewered on the \remcom{man's} \remcom{pole}.\\\hline
\end{remglyph}

\begin{prons}
\pronsrtwocone & \textbf{イ (I)} & --- \\\hline
\pronsrthreecone & --- & --- \\\hline
\end{prons}

\begin{remprons}
\hooglig{With} \comon{indication}.\\\hline
\end{remprons}

%{\middel}
% &  \mbox{()}\\\hline
\begin{somme}
六十以上 & sixty + with + upper = \textit{more than sixty} & ロク{\middel}ジュウ{\middel}イ{\middel}ジョウ \mbox{(ROKU{\middel}J\={U}{\middel}I{\middel}J\={O})}\\\hline
十度以内 & ten + degree + with + inside = \textit{within ten degrees} & ジュウ{\middel}ド{\middel}イ{\middel}ナイ \mbox{(J\={U}{\middel}DO{\middel}I{\middel}NAI)}\\\hline
私以外 & private (I) + with + outside = \textit{except for me} & わたし{\middel}イ{\middel}ガイ \mbox{(watashi{\middel}I{\middel}GAI)}\\\hline
\end{somme}

%\begin{decomp}{}
%\\\hline
%\\\hline
%\\\hline
%\multicolumn{}{|r|}{\textit{}}\\\hline
%\end{decomp}

\end{multicols}
\vhrulefill{2pt}
%%--------------------------------------------------------------------
%6--------------------------------------------------------------------
\begin{multicols}{2}
\begin{glyph}{Single (単)}{single}\label{glyph:single}
Single, simple.
\end{glyph}

\begin{comps}
\comprtwocone & 小 + 田 + 十 \\\hline
\comprthreecone & Small + Field + Ten \\\hline
\end{comps}

\begin{remglyph}
\hooglig{Single}, \remcom{small} \remcom{cross} in the \remcom{field}.\\\hline
\end{remglyph}

\begin{prons}
\pronsrtwocone & \textbf{タン (TAN)} & --- \\\hline
\pronsrthreecone & --- & --- \\\hline
\end{prons}

\begin{remprons}
\hooglig{Simple} \comon{ton} of bricks.\\\hline
\end{remprons}

%{\middel}
% &  \mbox{()}\\\hline
\begin{somme}
単語 & single + words = \textit{word} & タン{\middel}ゴ \mbox{(TAN{\middel}GO)}\\\hline
単数 & single + number = \textit{singular (opposite of plural)} & タン{\middel}スウ \mbox{(TAN{\middel}S\={U})}\\\hline
単刀直入 & single + sword + straight + enter = \textit{straight to the point; frank} & タン{\middel}トウ{\middel}チョク{\middel}ニュウ \mbox{(TAN{\middel}T\={O}{\middel}CH\={O}{\middel}NY\={U})}\\\hline
\end{somme}

%\begin{decomp}{}
%\\\hline
%\\\hline
%\\\hline
%\multicolumn{}{|r|}{\textit{}}\\\hline
%\end{decomp}

\end{multicols}
\vspace*{\fill}
\pagebreak
%%--------------------------------------------------------------------
%7--------------------------------------------------------------------
\vspace*{\fill}
\begin{multicols}{2}
\begin{glyph}{Versus (対)}{versus}\label{glyph:versus}
Versus.  The general sense is of things opposing each other in a confrontational way, or simply facing each other in a more neutral manner.
\end{glyph}

\begin{comps}
\comprtwocone & 文 + 寸 \\\hline
\comprthreecone & Word/Literature + Inch \\\hline
\end{comps}

\begin{remglyph}
\hooglig{Versus}, just like \remcom{inch}, are words that can be found in the \remcom{literature}.\\\hline
\end{remglyph}

\begin{prons}
\pronsrtwocone & \textbf{タイ (TAI)} & --- \\\hline
\pronsrthreecone & ツイ (TSUI) & --- \\\hline
\end{prons}

\begin{remprons}
\hooglig{Versus} \ucomon{Tsui-Goab} means it's going to be a \comon{tight} match.\\\hline
\textit{`Tsui-Goab'} is a mythological god of creation.\\\hline
\end{remprons}

%{\middel}
% &  \mbox{()}\\\hline
\begin{somme}
対立 & versus + stand = \textit{opposition; confrontation} & タイ{\middel}リツ \mbox{(TAI{\middel}RITSU)}\\\hline
対比 & versus + compare = \textit{comparison; contrast} & タイ{\middel}ヒ \mbox{(TAI{\middel}HI)}\\\hline
反対 & against + versus = \textit{opposite; contrary} & ハン{\middel}タイ \mbox{(HAN{\middel}TAI)}\\\hline
\end{somme}

%\begin{decomp}{}
%\\\hline
%\\\hline
%\\\hline
%\multicolumn{}{|r|}{\textit{}}\\\hline
%\end{decomp}

\end{multicols}
\vhrulefill{2pt}
%%--------------------------------------------------------------------
%8--------------------------------------------------------------------
\begin{multicols}{2}
\begin{glyph}{Meal (飯)}{meal}\label{glyph:meal}
Meal, (cooked) rice.
\end{glyph}

\begin{comps}
\comprtwocone & 食 + 反 \\\hline
\comprthreecone & Eat/Food + Against (\ref{glyph:against}) \\\hline
\end{comps}

\begin{remglyph}
\hooglig{The meal} we \remcom{ate} was \remcom{against} all standards.\\\hline
\end{remglyph}

\begin{prons}
\pronsrtwocone & \textbf{ハン (HAN)} & \textbf{めし (meshi)}\\\hline
\pronsrthreecone & --- & --- \\\hline
\end{prons}

\begin{remprons}
\hooglig{The meal} did not \comkun{mesh} well with the \comon{hunchback}.
\\\hline
\end{remprons}

%{\middel}
% &  \mbox{()}\\\hline
\begin{somme}
\rye{3}{Note that ご in the second example can also be written with the honorific prefix 御{\notcov}.}
飯 & \textit{meal; (cooked) rice} & めし \mbox{(meshi)}\\\hline
ご飯 & \textit{meal; (cooked) rice} & ご{\middel}ハン \mbox{(go{\middel}HAN)}\\\hline
\end{somme}

%\begin{decomp}{}
%\\\hline
%\\\hline
%\\\hline
%\multicolumn{}{|r|}{\textit{}}\\\hline
%\end{decomp}

\end{multicols}
\vspace*{\fill}
\pagebreak
%%--------------------------------------------------------------------
%9--------------------------------------------------------------------
\vspace*{\fill}
\begin{multicols}{2}
\begin{glyph}{Title (号)}{title}\label{glyph:title}
Title, sign.\\\hline
This is a versatile 漢字 that can refer to everything from traffic signals to newspaper editions.  The idea of a title or sign, however, is usually present in some way.
\end{glyph}

\begin{comps}
\comprtwocone & 口 + 勹 + 一 \\\hline
\comprthreecone & Mouth + Wrap/Embrace + One \\\hline
\end{comps}

\begin{remglyph}
\hooglig{The title} \remcom{embraced} by \remcom{one} \remcom{vampire}.\\\hline
\end{remglyph}

\begin{prons}
\pronsrtwocone & \textbf{ゴウ (G\={O})} & --- \\\hline
\pronsrthreecone & --- & --- \\\hline
\end{prons}

\begin{remprons}
\hooglig{Title} that is \comon{gory}.\\\hline
\end{remprons}

%{\middel}
% &  \mbox{()}\\\hline
\begin{somme}
信号 & trust + title = \textit{signal} & シン{\middel}ゴウ \mbox{(SHIN{\middel}G\={O})}\\\hline
番号 & order + title = \textit{number} & バン{\middel}ゴウ \mbox{(BAN{\middel}G\={O})}\\\hline
年号 & year + title = \textit{name of an era} & ネン{\middel}ゴウ \mbox{(NEN{\middel}G\={O})}\\\hline
\end{somme}

%\begin{decomp}{}
%\\\hline
%\\\hline
%\\\hline
%\multicolumn{}{|r|}{\textit{}}\\\hline
%\end{decomp}

\end{multicols}
\vhrulefill{2pt}
%%--------------------------------------------------------------------
%10-------------------------------------------------------------------
\begin{multicols}{2}
\begin{glyph}{Graduate (卒)}{graduate}\label{glyph:graduate}
Graduate, come to an end.\\\hline
A secondary meaning relates to soldiers.
\end{glyph}

\begin{comps}
\comprtwocone & 亠 + 人 + 十 \\\hline
\comprthreecone & Lid/Top + Man/Person + Ten \\\hline
\end{comps}

\begin{remglyph}
\hooglig{Graduation} of the \remcom{top} \remcom{10} \remcom{people}.\\\hline
\end{remglyph}

\begin{prons}
\pronsrtwocone & \textbf{ソツ (SOTSU)} & --- \\\hline
\pronsrthreecone & --- & --- \\\hline
\end{prons}

\begin{remprons}
\hooglig{Graduating} \comon{sots} on their graduation day.\\\hline
\end{remprons}

%{\middel}
% &  \mbox{()}\\\hline
\begin{somme}
\rye{3}{The second compound is actually a shortening of 大学卒業者 (ダイ{\middel}ガク{\middel}ソツ{\middel}ギョウ{\middel}シャ/DAI{\middel}GAKU{\middel}SOTSU{\middel}GY\={O}{\middel}SHA), from `university', `graduation', and `individual'.  This type of word creation is a common feature of Japanese.}
卒業 & graduate + business = \textit{graduation} & ソツ{\middel}ギョウ \mbox{(SOTSU{\middel}GY\={O})}\\\hline
大卒 & large + graduate = \textit{university graduate} & ダイ{\middel}ソツ \mbox{(DAI{\middel}SOTSU)}\\\hline
\end{somme}

%\begin{decomp}{}
%\\\hline
%\\\hline
%\\\hline
%\multicolumn{}{|r|}{\textit{}}\\\hline
%\end{decomp}

\end{multicols}
\vspace*{\fill}
\pagebreak
%%--------------------------------------------------------------------
%11-------------------------------------------------------------------
\vspace*{\fill}
\begin{multicols}{2}
\begin{glyph}{Sheet (枚)}{sheet}\label{glyph:sheet}
Sheet.\\\hline
This is the counting word that Japanese uses for flat, thin objects.
\end{glyph}

\begin{comps}
\comprtwocone & 木 + 攵 \\\hline
\comprthreecone & Tree + Whip \\\hline
\end{comps}

\begin{remglyph}
\hooglig{Sheets} are made by \remcom{whipping} up some \remcom{wood} concoction.\\\hline
\end{remglyph}

\begin{prons}
\pronsrtwocone & \textbf{マイ (MAI)} & --- \\\hline
\pronsrthreecone & --- & --- \\\hline
\end{prons}

\begin{remprons}
\hooglig{Sheet} contents \comon{migrating} to your brain.\\\hline
\end{remprons}

%{\middel}
% &  \mbox{()}\\\hline
\begin{somme}
百枚 & hundred + sheet = \textit{one hundred sheets} & ヒャク{\middel}マイ \mbox{(HYAKU{\middel}MAI)}\\\hline
枚数 & sheet + number = \textit{number of sheets} & マイ{\middel}スウ \mbox{(MAI{\middel}S\={U})}\\\hline
\end{somme}

%\begin{decomp}{}
%\\\hline
%\\\hline
%\\\hline
%\multicolumn{}{|r|}{\textit{}}\\\hline
%\end{decomp}

\end{multicols}
\vhrulefill{2pt}
%%--------------------------------------------------------------------
%12-------------------------------------------------------------------
\begin{multicols}{2}
\begin{glyph}{Hot (暑)}{hot}\label{glyph:hot}
Hot.\\\hline
This 漢字 relates solely to air temperature and weather; it's opposite is thus 寒, `cold' (Entry \ref{glyph:cold}).
\end{glyph}

\begin{comps}
\comprtwocone & 日 + 耂 \\\hline
\comprthreecone & Sun + Old \\\hline
\end{comps}

\begin{remglyph}
\hooglig{Hot}, \remcom{old} \remcom{sun}.\\\hline
\end{remglyph}

\begin{prons}
\pronsrtwocone & \textbf{ショ (SHO)} & \textbf{あつ (atsu)}\\\hline
\pronsrthreecone & --- & --- \\\hline
\end{prons}

\begin{remprons}
Hot atsumori play souvenirs available at the shop.\\\hline
%\rye{3}{\textit{`Atsumori'} is a Japanese Noh play.}
\end{remprons}

%{\middel}
% &  \mbox{()}\\\hline
\begin{somme}
暑い & \textit{hot (weather)} & あつ{\middel}い \mbox{(atsu{\middel}i)}\\\hline
暑がる & \textit{suffer from the heat} & あつ{\middel}がる \mbox{(atsu{\middel}garu)}\\\hline
暑中 & hot + middle = \textit{height of summer; hot season} & ショ{\middel}チュウ \mbox{(SHO{\middel}CH\={U})}\\\hline
\end{somme}

%\begin{decomp}{}
%\\\hline
%\\\hline
%\\\hline
%\multicolumn{}{|r|}{\textit{}}\\\hline
%\end{decomp}

\end{multicols}
\vspace*{\fill}
\pagebreak
%%--------------------------------------------------------------------
%13-------------------------------------------------------------------
\vspace*{\fill}
\begin{multicols}{2}
\begin{glyph}{Ride (乗)}{ride}\label{glyph:ride}
Ride.\\\hline
A secondary meaning relates to deceiving someone or being deceived oneself.
\end{glyph}

\begin{comps}
\comprtwocone & 禾 + 一 + 丨 \\\hline
\comprthreecone & Two-branch tree + One + Pole \\\hline
\end{comps}

\begin{remglyph}
\hooglig{A ride} for \remcom{one} \remcom{pole}-dancing a \remcom{two-branch tree}.\\\hline
\end{remglyph}

\begin{prons}
\pronsrtwocone & \textbf{ジョウ (J\={O})} & \textbf{の (no)}\\\hline
\rye{3}{The 訓読み can incorporate its ひらがな ending in certain compounds.}
\pronsrthreecone & --- & --- \\\hline
\end{prons}

\begin{remprons}
\hooglig{A ride} so fast it will rip off a \comkun{nob's} \comon{jaw}.\\\hline
\end{remprons}

%{\middel}
% &  \mbox{()}\\\hline
\begin{somme}
乗る & \textit{ride; get onto} & の{\middel}る \mbox{(no{\middel}ru)}\\\hline
乗り場 & ride + place = \textit{(bus) stop; place to get on} & の{\middel}り　ば \mbox{(no{\middel}ri ba)}\\\hline
同乗 & same + ride = \textit{ride together; share a ride} & ドウ{\middel}ジョウ \mbox{(D\={O}{\middel}J\={O})}\\\hline
\end{somme}

%\begin{decomp}{}
%\\\hline
%\\\hline
%\\\hline
%\multicolumn{}{|r|}{\textit{}}\\\hline
%\end{decomp}

\end{multicols}
\vhrulefill{2pt}
%%--------------------------------------------------------------------
%14-------------------------------------------------------------------
\begin{multicols}{2}
\begin{glyph}{Lady (婦)}{lady}\label{glyph:lady}
Lady, woman.\\\hline
It's worthwhile to see how this 漢字 differs visually from 寝 `sleep' (Entry \ref{glyph:sleep}), and 帰 `homeward' (Entry \ref{glyph:homeward}).
\end{glyph}

\begin{comps}
\comprtwocone & 女 + 彐彑 + 冖 + 巾 \\\hline
\comprthreecone & Woman + Pig head/snout + Cover/Crown + Cloth \\\hline
\end{comps}

\begin{remglyph}
\hooglig{Lady} \remcom{pig-head} \remcom{covered} in \remcom{woman's} \remcom{cloth}.\\\hline
\end{remglyph}

\begin{prons}
\pronsrtwocone & \textbf{フ (FU)} & --- \\\hline
\pronsrthreecone & --- & --- \\\hline
\end{prons}

\begin{remprons}
\hooglig{A lady's} \comon{foofy} tricks.\\\hline
\end{remprons}

%{\middel}
% &  \mbox{()}\\\hline
\begin{somme}
\rye{3}{Recall from Entry \ref{glyph:husband} that the final compound is one of the only instances in which 夫 is read with its less common 音読み of フウ.}
婦人 & lady + person = \textit{lady; woman} & フ{\middel}ジン \mbox{(FU{\middel}JIN)}\\\hline
婦人用 & lady + person + utilize = \textit{for ladies; women's} & フ{\middel}ジン{\middel}ヨウ \mbox{(FU{\middel}JIN{\middel}Y\={O})}\\\hline
夫婦 & husband + lady = \textit{husband and wife} & フウ{\middel}フ \mbox{(F\={U}{\middel}FU)}\\\hline
\end{somme}

%\begin{decomp}{}
%\\\hline
%\\\hline
%\\\hline
%\multicolumn{}{|r|}{\textit{}}\\\hline
%\end{decomp}

\end{multicols}
\vspace*{\fill}
\pagebreak
%%--------------------------------------------------------------------
%15-------------------------------------------------------------------
\vspace*{\fill}
\begin{multicols}{2}
\begin{glyph}{Nature (然)}{nature}\label{glyph:nature}
Nature.\\\hline
This 漢字 emphasizes the nature or personality of the character preceding it.
\end{glyph}

\begin{comps}
\comprtwocone & 月 + 犬 + 灬 \\\hline
\comprthreecone & Moon + Dog + Fire \\\hline
\end{comps}

\begin{remglyph}
\hooglig{Nature} of the \remcom{fire}-\remcom{hound} howling at the \remcom{moon}.\\\hline
\end{remglyph}

\begin{prons}
\pronsrtwocone & \textbf{ゼン (ZEN)} & --- \\\hline
\pronsrthreecone & ネン (NEN) & --- \\\hline
\rye{3}{You are likely to meet this reading in one commonly-used compound: 天然 (テン{\middel}ネン/TEN{\middel}NEN) `natural', from the 漢字 `heaven' and `nature'.}
\end{prons}

\begin{remprons}
\hooglig{The Nature} of \ucomon{Nenets'} \comon{Zen}.\\\hline
The \textit{`Nenets'} are an ethnic group native to northern arctic Russia.\\\hline
\end{remprons}

%{\middel}
% &  \mbox{()}\\\hline
\begin{somme}
\rye{3}{Note that the second compound features the only instance of 自 being read with its less common 音読み of シ.}
全然 & complete + nature = \textit{entirely; (not) at all} & ゼン{\middel}ゼン \mbox{(ZEN{\middel}ZEN)}\\\hline
自然 & self + nature = \textit{nature} & シ{\middel}ゼン \mbox{(SHI{\middel}ZEN)}\\\hline
\end{somme}

%\begin{decomp}{}
%\\\hline
%\\\hline
%\\\hline
%\multicolumn{}{|r|}{\textit{}}\\\hline
%\end{decomp}

\end{multicols}
\vhrulefill{2pt}
%%--------------------------------------------------------------------
%16-------------------------------------------------------------------
\begin{multicols}{2}
\begin{glyph}{Tone (調)}{tone}\label{glyph:tone}
Tone, tune.\\\hline
As an extension to toning and tuning comes the ideas of investigation, checking and preparation of things.
\end{glyph}

\begin{comps}
\comprtwocone & 言 + 周 \\\hline
\comprthreecone & Say/Speech + Around (\ref{glyph:around}) \\\hline
\end{comps}

\begin{remglyph}
\hooglig{Tune} your \remcom{sayings} to be centred \remcom{around} the Bible.\\\hline
\end{remglyph}

\begin{prons}
\pronsrtwocone & \textbf{チョウ (CH\={O})} & \textbf{しら (shira)}\\\hline
\pronsrthreecone & --- & ととの (totono) \\\hline
\end{prons}

\begin{remprons}
\hooglig{Investigating} is such a \comon{chore} and can result in \comkun{sheer ugliness} when \ucomkun{tot honours} are involved.\\\hline
\textit{`Tot honours'} are a small amount of a strong alcoholic drink such as whiskey/brandy.\\\hline
\end{remprons}

%{\middel}
% &  \mbox{()}\\\hline
\begin{somme}
調べる & \textit{investigate; inspect} & しら{\middel}べる \mbox{(shira{\middel}beru)}\\\hline
強調 & strong + tone = \textit{emphasis; stress} & キョウ{\middel}チョウ \mbox{(KY\={O}{\middel}CH\={O})}\\\hline
単調 & single + tone = \textit{monotonous; dull} & タン{\middel}チョウ \mbox{(TAN{\middel}CH\={O})}\\\hline
\end{somme}

%\begin{decomp}{}
%\\\hline
%\\\hline
%\\\hline
%\multicolumn{}{|r|}{\textit{}}\\\hline
%\end{decomp}

\end{multicols}
\vspace*{\fill}
\pagebreak
%%--------------------------------------------------------------------
%17-------------------------------------------------------------------
\vspace*{\fill}
\begin{multicols}{2}
\begin{glyph}{Keep (留)}{keep}\label{glyph:keep}
Keep, detain.
\end{glyph}

\begin{comps}
\comprtwocone & 厶 + 刀 + 田 + 丶 \\\hline
\comprthreecone & Self/Private + Sword/Knife + Rice Field + Dot \\\hline
\end{comps}

\begin{remglyph}
\hooglig{Keep} the \remcom{private} \remcom{jelly-beans} picked up in the \remcom{rice-field}, even if it means protecting it with a \remcom{sword}.\\\hline
\end{remglyph}

\begin{prons}
\pronsrtwocone & \textbf{リュウ (RY\={U})} & \textbf{と (to)}\\\hline
\pronsrthreecone & ル (RU) & --- \\\hline
\end{prons}

\begin{remprons}
\hooglig{Keep} the \comkun{torrent} of \ucomon{rookies} afflicted with \comon{leukaemia}.\\\hline
\end{remprons}

%{\middel}
% &  \mbox{()}\\\hline
\begin{somme}
\rye{3}{The verbs below have some overlap with an identically pronounced pair we learned back in Entry \ref{glyph:stop} (止まる / 止める).  The two here, however, have more of a sense of keeping in place of being kept in place.}
留まる (intr) & \textit{to be kept in place} & と{\middel}まる \mbox{(to{\middel}maru)}\\\hline
留める (tr) & \textit{keep in place; detain} & と{\middel}める \mbox{(to{\middel}meru)}\\\hline
留置 & keep + put = \textit{detention; custody} & リュウ{\middel}チ \mbox{(RY\={U}{\middel}CHI)}\\\hline
\end{somme}

%\begin{decomp}{}
%\\\hline
%\\\hline
%\\\hline
%\multicolumn{}{|r|}{\textit{}}\\\hline
%\end{decomp}

\end{multicols}
\vhrulefill{2pt}
%%--------------------------------------------------------------------
%18-------------------------------------------------------------------
\begin{multicols}{2}
\begin{glyph}{Heat (熱)}{heat}\label{glyph:heat}
Heat, hot (in aspects other than weather), and the ideas of enthusiasm and passion.  The opposite of this 漢字 is 冷, `cool' (Entry \ref{glyph:cool}).
\end{glyph}

\begin{comps}
\comprtwocone & 土 + 儿 + 灬 + 丸 \\\hline
\comprthreecone & Earth + Man/Legs + Fire + Round (\ref{glyph:round})\\\hline
\end{comps}

\begin{remglyph}
\hooglig{Heat} from the \remcom{round} \remcom{fireballs} in the centre of the \remcom{Earth}.\\\hline
\end{remglyph}

\begin{prons}
\pronsrtwocone & \textbf{ネツ (NETSU)} & \textbf{あつ (atsu)}\\\hline
\pronsrthreecone & --- & --- \\\hline
\end{prons}

\begin{remprons}
\hooglig{Passionate} \comon{network} of love in the \comkun{Atsumori} play.\\\hline
\end{remprons}

%{\middel}
% &  \mbox{()}\\\hline
\begin{somme}
熱 & \textit{heat; a fever} & ネツ \mbox{(NETSU)}\\\hline
熱い & \textit{hot (to the touch); heated} & あつ{\middel}い \mbox{(atsu{\middel}i)}\\\hline
熱愛 & heat + love = \textit{passionate love; devotion} & ネツ{\middel}アイ \mbox{(NETSU{\middel}AI)}\\\hline
\end{somme}

%\begin{decomp}{}
%\\\hline
%\\\hline
%\\\hline
%\multicolumn{}{|r|}{\textit{}}\\\hline
%\end{decomp}

\end{multicols}
\vspace*{\fill}
\pagebreak
%%--------------------------------------------------------------------
%19-------------------------------------------------------------------
\vspace*{\fill}
\begin{multicols}{2}
\begin{glyph}{Use (使)}{use}\label{glyph:use}
Use.\\\hline
This is the last of a group of five 漢字 that we can now review as a whole.  The building patter for the present character has been as follows: 史 `history' (Entry \ref{glyph:history}), to 吏 `officer' (Entry \ref{glyph:officer}), to 使.  A pair of similar-looking but unrelated 漢字 were constructed in this way: 更 `anew' (Entry \ref{glyph:anew}), to 便 `delivery' (Entry \ref{glyph:delivery.})
\end{glyph}

\begin{comps}
\comprtwocone & 亻 + 吏 \\\hline
\comprthreecone & Man/Person + Officer (\ref{glyph:officer}) \\\hline
\end{comps}

\begin{remglyph}
\hooglig{Useful} \remcom{officer} \remcom{man}.\\\hline
\end{remglyph}

\begin{prons}
\pronsrtwocone & \textbf{シ (SHI)} & \textbf{つか (tsuka)}\\\hline
\rye{3}{The 訓読み can become voiced in the second position, although never when forming verbs.}
\pronsrthreecone & --- & --- \\\hline
\end{prons}

\begin{remprons}
\hooglig{Use} the \comon{shield} and the \comkun{tsuka} together.\\\hline
A \textit{`tsuka'} is a Katana's handle.\\\hline
\end{remprons}

%{\middel}
% &  \mbox{()}\\\hline
\begin{somme}
使う & \textit{use} & つか{\middel}う \mbox{(tsuka{\middel}u)}\\\hline
使用 & use + utilize = \textit{use; utilize} & シン{\middel}ヨウ \mbox{(SHIN{\middel}Y\={O})}\\\hline
天使 & heaven + use = \textit{angel} & テン{\middel}シ \mbox{(TEN{\middel}SHI)}\\\hline
\end{somme}

%\begin{decomp}{}
%\\\hline
%\\\hline
%\\\hline
%\multicolumn{}{|r|}{\textit{}}\\\hline
%\end{decomp}

\end{multicols}
\vhrulefill{2pt}
%%--------------------------------------------------------------------
%20-------------------------------------------------------------------
\begin{multicols}{2}
\begin{glyph}{Remain (残)}{remain}\label{glyph:remain}
Remain.\\\hline
In an interesting connection, a secondary meaning relates to cruelty.
\end{glyph}

\begin{comps}
\comprtwocone & 歹 + 戈 + 二 \\\hline
\comprthreecone & Death/Decay + Spear/Halberd + Two \\\hline
\end{comps}

\begin{remglyph}
\hooglig{The remaining} \remcom{two} corpses were \remcom{speared} and left to \remcom{decay}.\\\hline
\end{remglyph}

\begin{prons}
\pronsrtwocone & \textbf{ザン (ZAN)} & \textbf{のこ (noko)}\\\hline
\pronsrthreecone & --- & --- \\\hline
\end{prons}

\begin{remprons}
\hooglig{Remaining} \comon{Zangetsu} needs to \comkun{knock off}.\\\hline
The \textit{`Zangetsu'} from Bleach.\\\hline
\end{remprons}

%{\middel}
% &  \mbox{()}\\\hline  
\begin{somme}
\rye{3}{Note the mix of 音読み and 訓読み in the final compound (and the voicing of 高).}
残る & \textit{remain; be left} & のこ{\middel}る \mbox{(noko{\middel}ru)}\\\hline
残業 & remain + business = \textit{overtime work} & ザン{\middel}ギョウ \mbox{(ZAN{\middel}GY\={O})}\\\hline
残高 & remain + tall = \textit{(account) balance} & ザン{\middel}だか \mbox{(ZAN{\middel}daka)}\\\hline
\end{somme}

%\begin{decomp}{}
%\\\hline
%\\\hline
%\\\hline
%\multicolumn{}{|r|}{\textit{}}\\\hline
%\end{decomp}

\end{multicols}
\vspace*{\fill}
\pagebreak
%%--------------------------------------------------------------------
%21-------------------------------------------------------------------
\vspace*{\fill}
\begin{multicols}{2}
\begin{glyph}{Quiet (静)}{quiet}\label{glyph:quiet}
Quiet.
\end{glyph}

\begin{comps}
\comprtwocone & 青 + 争 \\\hline
\comprthreecone & Blue + Dispute (\ref{glyph:dispute}) \\\hline
\end{comps}

\begin{remglyph}
\hooglig{A quiet} \remcom{dispute} between people feeling \remcom{blue}.\\\hline
\end{remglyph}

\begin{prons}
\pronsrtwocone & \textbf{セイ (SEI)} & \textbf{しず (shizu)}\\\hline
\pronsrthreecone & ジョウ (J\={O}) & --- \\\hline
\end{prons}

\begin{remprons}
\hooglig{Quietly} \comkun{she zoomed} into his \ucomon{jaw} while \comon{sailing}.\\\hline
\end{remprons}

%{\middel}
% &  \mbox{()}\\\hline
\begin{somme}
静か & \textit{quiet; peaceful} & しず{\middel}か \mbox{(shizu{\middel}ka)}\\\hline
安静 & ease + quiet = \textit{rest; repose} & アン{\middel}セイ \mbox{(AN{\middel}SEI)}\\\hline
静物画 & quiet + thing + painting = \textit{still life} & セイ{\middel}ブツ{\middel}ガ \mbox{(SEI{\middel}BUTSU{\middel}GA)}\\\hline
\end{somme}

%\begin{decomp}{}
%\\\hline
%\\\hline
%\\\hline
%\multicolumn{}{|r|}{\textit{}}\\\hline
%\end{decomp}

\end{multicols}
\vhrulefill{2pt}
%%--------------------------------------------------------------------
%22-------------------------------------------------------------------
\begin{multicols}{2}
\begin{glyph}{Gift (贈)}{gift}\label{glyph:gift}
Gift, present.
\end{glyph}

\begin{comps}
\comprtwocone & 貝 + 艹 + 田 + 日 \\\hline
\comprthreecone & Shellfish + Grass + Rice Field + Sun \\\hline
\end{comps}

\begin{remglyph}
\hooglig{The gift} of \remcom{shellfish} wrapped in \remcom{herbs} and found \remcom{clearly} in the \remcom{field}.\\\hline
\end{remglyph}

\begin{prons}
\pronsrtwocone & --- & \textbf{おく (oku)}\\\hline
\pronsrthreecone & ゾウ、ソウ (Z\={O}, S\={O}) & --- \\\hline
\end{prons}

\begin{remprons}
\hooglig{The gift} of being \comkun{occupied} with \ucomon{zorbing} until you feel \ucomon{sore}.\\\hline
\end{remprons}

%{\middel}
% &  \mbox{()}\\\hline
\begin{somme}
贈る & \textit{present (a gift)} & おく{\middel}る \mbox{(oku{\middel}ru)}\\\hline
贈り物 & gift + thing = \textit{gift; present} & おく{\middel}り　もの \mbox{(oku{\middel}ri mono)}\\\hline
\end{somme}

%\begin{decomp}{}
%\\\hline
%\\\hline
%\\\hline
%\multicolumn{}{|r|}{\textit{}}\\\hline
%\end{decomp}

\end{multicols}
\vspace*{\fill}
\pagebreak
%%--------------------------------------------------------------------
%23-------------------------------------------------------------------
\vspace*{\fill}
\begin{multicols}{2}
\begin{glyph}{Vegetable (菜)}{vegetable}\label{glyph:vegetable}
Vegetable.
\end{glyph}

\begin{comps}
\comprtwocone & 艹 + 采 \\\hline
\comprthreecone & Grass/Herb/Plant + Appearance (\ref{glyph:appearance}) \\\hline
\end{comps}

\begin{remglyph}
\hooglig{Vegetable} has some \remcom{plant}-like \remcom{appearance}.\\\hline
\end{remglyph}

\begin{prons}
\pronsrtwocone & \textbf{サイ (SAI)} & --- \\\hline
\pronsrthreecone & --- & な (na) \\\hline
\end{prons}

\begin{remprons}
\hooglig{Vegetables} to \comon{psyche} you up and remove the \ucomkun{numbness}.\\\hline
\end{remprons}

%{\middel}
% &  \mbox{()}\\\hline
\begin{somme}
野菜 & field + vegetable = \textit{vegetables} & ヤ{\middel}サイ \mbox{(YA{\middel}SAI)}\\\hline
菜食 & vegetable + eat = \textit{vegetarian diet} & サイ{\middel}ショク \mbox{(SAI{\middel}SHOKU)}\\\hline
菜園 & vegetable + garden = \textit{vegetable garden} & サイ{\middel}エン \mbox{(SAI{\middel}EN)}\\\hline
\end{somme}

%\begin{decomp}{}
%\\\hline
%\\\hline
%\\\hline
%\multicolumn{}{|r|}{\textit{}}\\\hline
%\end{decomp}

\end{multicols}
\vhrulefill{2pt}
%%--------------------------------------------------------------------
%24-------------------------------------------------------------------
\begin{multicols}{2}
\begin{glyph}{Miscellaneous (雑)}{miscellaneous}\label{glyph:miscellaneous}
Miscellaneous.
\end{glyph}

\begin{comps}
\comprtwocone & 九 + 木 + 隹 \\\hline
\comprthreecone & Nine (\ref{glyph:nine}) + Tree + Old Bird \\\hline
\end{comps}

\begin{remglyph}
\hooglig{Miscellaneous} \remcom{small birds} in the \remcom{tree} of up to \remcom{9} kinds.\\\hline
\end{remglyph}

\begin{prons}
\pronsrtwocone & \textbf{ザツ (ZATSU)} & --- \\\hline
\pronsrthreecone & ゾウ (Z\={O}) & --- \\\hline
\end{prons}

\begin{remprons}
\hooglig{Miscellaneous} \comon{zuts} were uttered during the \ucomon{zorbing} experience.\\\hline
\textit{`Zut'} means ``Damn it'' in French.\\\hline
\end{remprons}

%{\middel}
% &  \mbox{()}\\\hline
\begin{somme}
雑音 & miscellaneous + sound = \textit{noise; interference} & ザツ{\middel}オン \mbox{(ZATSU{\middel}ON)}\\\hline
雑事 & miscellaneous + matter = \textit{miscellaneous affairs} & ザツ{\middel}ジ \mbox{(ZATSU{\middel}JI)}\\\hline
雑用 & miscellaneous + utilize = \textit{miscellaneous things to do} & ザツ{\middel}ヨウ \mbox{(ZATSU{\middel}Y\={O})}\\\hline
\end{somme}

%\begin{decomp}{}
%\\\hline
%\\\hline
%\\\hline
%\multicolumn{}{|r|}{\textit{}}\\\hline
%\end{decomp}

\end{multicols}
\vspace*{\fill}
\pagebreak
%%--------------------------------------------------------------------
%25-------------------------------------------------------------------
\vspace*{\fill}
\begin{multicols}{2}
\begin{glyph}{Without (無)}{without}\label{glyph:without}
This is another important negating 漢字, along with 不 `not' (Entry \ref{glyph:not}), and 非 `un-' (Entry \ref{glyph:un}).  It imparts a sense of `-less' or `-free' in many of the compounds in which it appears.
\end{glyph}

\begin{comps}
\comprtwocone & 無 \\\hline
\comprthreecone & Nothingness \\\hline
\end{comps}

\begin{remglyph}
An uncommon 漢字 radical.\\\hline
\end{remglyph}

\begin{prons}
\pronsrtwocone & \textbf{ム、ブ (MU, BU)} & \textbf{な (na)}\\\hline
\rye{3}{ム is by far the more common of these readings, although ブ occurs in a few important compounds.\\\hline
Except when it appears on its own, な will be found only in the second or third position.}
\pronsrthreecone & --- & --- \\\hline
\end{prons}

\begin{remprons}
\hooglig{Without} \comon{muti} that \comkun{numbs} the pain, \comon{bulleying} would have been a terrible experience.\\\hline
\end{remprons}

%{\middel}
% &  \mbox{()}\\\hline
\begin{somme}
無い & \textit{there is no\ldots} & ない \mbox{(na{\middel}i)}\\\hline
無料 & without + fee = \textit{free of charge} & ム{\middel}リョウ \mbox{(MU{\middel}RY\={O})}\\\hline
無事 & without + matter = \textit{without incident; safe and sound} & ブ{\middel}ジ \mbox{(BU{\middel}JI)}\\\hline
\end{somme}

%\begin{decomp}{}
%\\\hline
%\\\hline
%\\\hline
%\multicolumn{}{|r|}{\textit{}}\\\hline
%\end{decomp}

\end{multicols}
\vhrulefill{2pt}
%%--------------------------------------------------------------------
%26-------------------------------------------------------------------
\begin{multicols}{2}
\begin{glyph}{Yard (庭)}{yard}\label{glyph:yard}
Yard, garden.
\end{glyph}

\begin{comps}
\comprtwocone & 广 + 廷 \\\hline
\comprthreecone & Slanting roof + Court (\ref{glyph:court}) \\\hline
\end{comps}

\begin{remglyph}
\hooglig{The yard} with the \remcom{slanting roof} is a beautiful \remcom{courtyard}.\\\hline
\end{remglyph}

\begin{prons}
\pronsrtwocone & \textbf{テイ (TEI)} & \textbf{にわ (niwa)}\\\hline
\pronsrthreecone & --- & --- \\\hline
\end{prons}

\begin{remprons}
\hooglig{The yard} had a \comkun{knee-walker} \comon{taped} to a tree.\\\hline
\end{remprons}

%{\middel}
% &  \mbox{()}\\\hline
\begin{somme}
庭 & \textit{yard; garden} & にわ \mbox{(niwa)}\\\hline
庭園 & yard + garden = \textit{garden} & テイ{\middel}エン \mbox{(TEI{\middel}EN)}\\\hline
家庭 & house + yard = \textit{home; household} & カ{\middel}テイ \mbox{(KA{\middel}TEI)}\\\hline
\end{somme}

%\begin{decomp}{}
%\\\hline
%\\\hline
%\\\hline
%\multicolumn{}{|r|}{\textit{}}\\\hline
%\end{decomp}

\end{multicols}
\vspace*{\fill}
\pagebreak
%%--------------------------------------------------------------------
%27-------------------------------------------------------------------
\vspace*{\fill}
\begin{multicols}{2}
\begin{glyph}{Reciprocal (相)}{reciprocal}\label{glyph:reciprocal}
Reciprocal, each other.\\\hline
An important secondary meaning is present when this 漢字 appears as the final character in certain compounds; in such instances it can indicate a government minister (as in the final example below).
\end{glyph}

\begin{comps}
\comprtwocone & 木 + 目 \\\hline
\comprthreecone & Tree + Eye \\\hline
\end{comps}

\begin{remglyph}
\hooglig{Each one} has a \remcom{wood} beam in their \remcom{eyes}.\\\hline
\end{remglyph}

\begin{prons}
\pronsrtwocone & \textbf{ソウ、ショウ (S\={O}, SH\={O})} & \textbf{あい (ai)}\\\hline
\rye{3}{ソウ is the most common reading here, with ショウ only used outside the first position to indicate government ministers.}
\pronsrthreecone & --- & --- \\\hline
\end{prons}

\begin{remprons}
\hooglig{Each other} \comon{shore} was an \comkun{icon} of \comon{soreness}.\\\hline
\end{remprons}

%{\middel}
% &  \mbox{()}\\\hline
\begin{somme}
相手 & reciprocal + hand = \textit{the other party; opponent} & あい{\middel}て \mbox{(ai{\middel}te)}\\\hline
相続 & reciprocal + continue = \textit{inheritance; succession} & ソウ{\middel}ゾク \mbox{(S\={O}{\middel}ZOKU)}\\\hline
首相 & neck (leader) + reciprocal (minister) = \textit{prime minister} & シュ{\middel}ショウ \mbox{(SHU{\middel}SH\={O})}\\\hline
\end{somme}

%\begin{decomp}{}
%\\\hline
%\\\hline
%\\\hline
%\multicolumn{}{|r|}{\textit{}}\\\hline
%\end{decomp}

\end{multicols}
\vhrulefill{2pt}
%%--------------------------------------------------------------------
%28-------------------------------------------------------------------
\begin{multicols}{2}
\begin{glyph}{Long (長)}{long}\label{glyph:long}
Long.\\\hline
This 漢字 can refer to the `chief' or `head' of entities such as cities and companies (the final compound provides an example).
\end{glyph}

\begin{comps}
\comprtwocone & 長 \\\hline
\comprthreecone & Long \\\hline
\end{comps}

\begin{remglyph}
Part of the 漢字 radicals (\#{168}).\\\hline
\end{remglyph}

\begin{prons}
\pronsrtwocone & \textbf{チョウ (CH\={O})} & \textbf{なが (naga)}\\\hline
\pronsrthreecone & --- & --- \\\hline
\end{prons}

\begin{remprons}
Long list of chores to be done in Nagasaki.\\\hline
\end{remprons}

%{\middel}
% &  \mbox{()}\\\hline
\begin{somme}
\rye{3}{Recall from Entry \ref{glyph:eternal} that the identically pronounced 永{\middel}い (なが{\middel}い/naga{\middel}i) refers to `long' in terms of time.}
長い & \textit{long} & なが{\middel}い \mbox{(naga{\middel}i)}\\\hline
長期 & long + period = \textit{long-term} & チョウ{\middel}キ \mbox{(CH\={O}{\middel}KI)}\\\hline
社長 & company + long (chief) = \textit{company president} & シャ{\middel}チョウ \mbox{(SHA{\middel}CH\={O})}\\\hline
\end{somme}

%\begin{decomp}{}
%\\\hline
%\\\hline
%\\\hline
%\multicolumn{}{|r|}{\textit{}}\\\hline
%\end{decomp}

\end{multicols}
\vspace*{\fill}
\pagebreak
%%--------------------------------------------------------------------
%29-------------------------------------------------------------------
\vspace*{\fill}
\begin{multicols}{2}
\begin{glyph}{Courtesy (礼)}{courtesy}\label{glyph:courtesy}
Courtesy, politeness.
\end{glyph}

\begin{comps}
\comprtwocone & 礻 + 乙乚 \\\hline
\comprthreecone & Altar/Display + Second/Latter \\\hline
\end{comps}

\begin{remglyph}
\hooglig{Courtesy} at the \remcom{alter} is \remcom{second} to nothing.\\\hline
\end{remglyph}

\begin{prons}
\pronsrtwocone & \textbf{レイ (REI)} & --- \\\hline
\pronsrthreecone & ライ (RAI) & --- \\\hline
\end{prons}

\begin{remprons}
\hooglig{Courtesy} of making a \ucomon{rhyming} poem about the \comon{rain}.\\\hline
\end{remprons}

%{\middel}
% &  \mbox{()}\\\hline
\begin{somme}
失礼 & lose + courtesy = \textit{discourtesy; rudeness} & シツ{\middel}レイ \mbox{(SHITSU{\middel}REI)}\\\hline
無礼 & without + courtesy = \textit{impoliteness; rudeness} & ブ{\middel}レイ \mbox{(BU{\middel}REI)}\\\hline
礼服 & courtesy + clothes = \textit{formal dress} & レイ{\middel}フク \mbox{(REI{\middel}FUKU)}\\\hline
\end{somme}

%\begin{decomp}{}
%\\\hline
%\\\hline
%\\\hline
%\multicolumn{}{|r|}{\textit{}}\\\hline
%\end{decomp}

\end{multicols}
\vhrulefill{2pt}
%%--------------------------------------------------------------------
%30-------------------------------------------------------------------
\begin{multicols}{2}
\begin{glyph}{Finish (終)}{finish}\label{glyph:finish}
Finish.
\end{glyph}

\begin{comps}
\comprtwocone & 糸 + 冬 \\\hline
\comprthreecone & Thread + Winter (\ref{glyph:winter}) \\\hline
\end{comps}

\begin{remglyph}
\hooglig{Finished} when your \remcom{thread} has reached the end of \remcom{winter}.\\\hline
\end{remglyph}

\begin{prons}
\pronsrtwocone & \textbf{シュウ (SH\={U})} & \textbf{お (o)}\\\hline
\pronsrthreecone & --- & --- \\\hline
\end{prons}

\begin{remprons}
\hooglig{When finished}, you are \comkun{obliged} to take off your \comon{shoes}.\\\hline
\end{remprons}

%{\middel}
% &  \mbox{()}\\\hline
\begin{somme}
終わる & \textit{finish; come to an end} & お{\middel}わる \mbox{(o{\middel}waru)}\\\hline
終結 & finish + bind = \textit{conclusion; end} & シュウ{\middel}ケツ \mbox{(SH\={U}{\middel}KETSU)}\\\hline
終点 & finish + point = \textit{last stop; end of the line} & シュウ{\middel}テン \mbox{(SH\={U}{\middel}TEN)}\\\hline
\end{somme}

%\begin{decomp}{}
%\\\hline
%\\\hline
%\\\hline
%\multicolumn{}{|r|}{\textit{}}\\\hline
%\end{decomp}

\end{multicols}
\vspace*{\fill}
\pagebreak
%%--------------------------------------------------------------------
%---------------------------------------------------------------------
